% Generated draft: requires mathematical and source review.
\documentclass[11pt]{article}
\usepackage[margin=25mm]{geometry}
\usepackage{fontspec}
\setmainfont{Latin Modern Roman}
\usepackage{amsmath,amssymb,mathtools,unicode-math}
\setmathfont{Latin Modern Math}
\usepackage{enumitem,needspace,xurl,hyperref}
\setlength{\parindent}{0pt}
\setlength{\parskip}{5pt}
\setlength{\emergencystretch}{3em}
\newcommand{\N}{\mathbb{N}}
\newcommand{\Z}{\mathbb{Z}}
\newcommand{\Q}{\mathbb{Q}}
\newcommand{\R}{\mathbb{R}}
\newcommand{\C}{\mathbb{C}}
\newcommand{\F}{\mathbb{F}}
\newcommand{\A}{\mathbb{A}}
\newcommand{\PP}{\mathbb{P}}
\newcommand{\E}{\mathbb{E}}
\newcommand{\eps}{\varepsilon}
\newcommand{\dd}{\,\mathrm{d}}
\begin{document}
\section*{UnsolvedMath ID: 701373}
\textbf{Source problem code:} catalogue entry
\subsection{Definitions and mathematical statement}
Let \(p\) be a prime. A group \(G\) is periodic if every element of \(G\) has finite order. A non-trivial \(p\)-element is an element of order \(p^k\) for some integer \(k\geq 1\). An automorphism \(\alpha\in\operatorname{Aut}(G)\) is fixed-point-free if
\[
\{g\in G:\alpha(g)=g\}=\{1_G\}.
\]
An automorphism \(\alpha\) of order \(p\) is called a splitting automorphism here if
\[
 g\,\alpha(g)\,\alpha^2(g)\cdots\alpha^{p-1}(g)=1_G
 \qquad\text{for every }g\in G,
\]
where the factors are multiplied in the displayed order.

\textbf{Part (b), the remaining question in the supplied record.} For every prime \(p\), every periodic group \(G\), and every fixed-point-free automorphism \(\alpha\in\operatorname{Aut}(G)\) of order exactly \(p\), does the hypothesis that \(G\) has no non-trivial \(p\)-element imply that \(\alpha\) is a splitting automorphism? Equivalently, under these hypotheses, must
\[
\forall g\in G,\qquad
 g\,\alpha(g)\cdots\alpha^{p-1}(g)=1_G
\]
hold? No finiteness assumption on \(G\) is imposed.

\textbf{Archived part (a).} The original accompanying question asks whether, for every prime \(p\), a periodic group admitting a fixed-point-free automorphism of order \(p\) necessarily has no non-trivial \(p\)-element. The supplied archive reports a negative answer: such a group may contain a non-trivial \(p\)-element. This is a supplied historical resolution attributed to E. B. Durakov and A. I. Sozutov, not an independently verified result in this draft. The supplied record labels the combined problem partially solved and retains part (b) as the remaining question.
\subsection{Short English statement}
Part (b) asks whether a fixed-point-free automorphism of prime order \(p\) of a periodic group with no nonidentity elements of \(p\)-power order must satisfy the splitting identity for every group element. The supplied archive reports that part (a), which asks whether the absence of such elements follows automatically, has a negative answer.
\subsection{Sources}
\begin{itemize}
\item Supplied local UnsolvedMath registry record, permanent ID \(701373\), exact source problem code \texttt{catalogue entry}, title \emph{Kourovka Notebook Problem 17.71}; upstream revision \texttt{d371ee917ac9647fbb06cd6fc529deec6f3dda82}.
\item V. D. Mazurov, proposer, 2010. \emph{The Kourovka Notebook (Unsolved Problems in Group Theory)}, 17th issue (2010), Problem 17.71; the supplied background identifies an update in the 21st issue (2026), arXiv:1401.0300v48. The supplied collection description names E. I. Khukhro and V. D. Mazurov as editors.
\item Supplied Archive of Solved Problems attribution for part (a): E. B. Durakov and A. I. Sozutov, \emph{Algebra Logic}, \textbf{52}, no. 5 (2013), 422--425. The supplied archive reports that a non-trivial \(p\)-element may occur. No article title or direct link is supplied.
\item The supplied status review is dated \texttt{2026-10-16}, labels the problem \texttt{partially\_solved}, records its evidence as \texttt{published}, and states that part (a) is archived as solved while the remaining part is in the main section. This is supplied dated curation, not an independently verified literature assessment.
\end{itemize}
\end{document}
